\documentclass{article}
\usepackage[utf8]{inputenc}
\usepackage{geometry}
\usepackage{graphicx}
\usepackage{hyperref}
\usepackage{amsmath}
\geometry{margin=1in}

\title{FAASI-CORE: A Benchmark Framework for Evaluating Autonomous Agent Reliability in Long-Horizon Tool-Augmented Workflows}
\author{David Carmel Alex \\ Fusion Civilization Research Institute (FCRI)}
\date{}

\begin{document}
\maketitle

\begin{abstract}
FAASI-CORE introduces a benchmark framework for evaluating autonomous agent reliability under long-horizon operational execution.
\end{abstract}

\section{Introduction}
Autonomous agents require benchmark methodologies beyond narrow completion metrics.

\section{Related Work}
Prior benchmark ecosystems include HELM, SWE-bench, AgentBench, and GAIA.

\section{Benchmark Taxonomy}
Task families include tool use, filesystem operations, recovery tasks, ambiguity resolution, and long-horizon workflows.

\section{Methodology}
FAASI evaluates task success, tool reliability, recovery intelligence, memory integrity, safety compliance, ambiguity governance, stability, and efficiency.

\section{Scoring Formalism}
\[
FAASI = \sum_{i=1}^{n} w_i m_i
\]

\section{Threats to Validity}
Prototype maturity, evaluation bias, and benchmark scope limitations remain active concerns.

\section{Ethics}
Safety-conscious benchmark design is required for autonomous system evaluation.

\bibliographystyle{plain}
\bibliography{references}

\end{document}
